% !TeX root = ../main.tex

\section{解答}

\begin{comment}
	\subsection{問題1}
	
	\begin{solution}
		% 解答を書く
	\end{solution}
	
	\begin{checknote}
		% 補足・検算・注意点を書く
	\end{checknote}
	
\end{comment}

\begin{proof}{A.6.(2)}
	(1) より，ある自然数 $N_0$ が存在して，$n\geq N_0$ ならば
	\[
	|b_n|\geq \frac{|\beta|}{2}
	\]
	が成り立つ。
	
	示したいことは
	\[
	\lim_{n\to\infty}\frac{a_n}{b_n}
	=
	\frac{\alpha}{\beta}
	\]
	である。すなわち，任意の $\epsilon>0$ に対して，ある自然数 $N$ が存在して，
	$n\geq N$ ならば
	\[
	\left|
	\frac{a_n}{b_n}-\frac{\alpha}{\beta}
	\right|<\epsilon
	\]
	となることを示せばよい。
	
	まず差を通分する。
	\[
	\frac{a_n}{b_n}-\frac{\alpha}{\beta}
	=
	\frac{\beta a_n-\alpha b_n}{\beta b_n}
	\]
	
	分子を次のように変形する。
	\[
	\beta a_n-\alpha b_n
	=
	\beta a_n-\beta\alpha+\beta\alpha-\alpha b_n
	\]
	したがって，
	\[
	\beta a_n-\alpha b_n
	=
	\beta(a_n-\alpha)+\alpha(\beta-b_n)
	\]
	である。
	
	よって，
	\[
	\left|
	\frac{a_n}{b_n}-\frac{\alpha}{\beta}
	\right|
	=
	\left|
	\frac{\beta(a_n-\alpha)+\alpha(\beta-b_n)}{\beta b_n}
	\right|
	\]
	である。
	三角不等式を用いると，
	\[
	\left|
	\frac{a_n}{b_n}-\frac{\alpha}{\beta}
	\right|
	\leq
	\frac{|\beta||a_n-\alpha|+|\alpha||\beta-b_n|}
	{|\beta||b_n|}
	\]
	となる。
	
	ここで，$n\geq N_0$ ならば
	\[
	|b_n|\geq \frac{|\beta|}{2}
	\]
	であるから，
	\[
	\frac{1}{|b_n|}
	\leq
	\frac{2}{|\beta|}
	\]
	が成り立つ。
	
	したがって，$n\geq N_0$ のとき
	\[
	\left|
	\frac{a_n}{b_n}-\frac{\alpha}{\beta}
	\right|
	\leq
	\frac{2}{|\beta|}|a_n-\alpha|
	+
	\frac{2|\alpha|}{|\beta|^2}|b_n-\beta|
	\]
	を得る。
	
	ここで，任意に $\epsilon>0$ を取る。
	
	$a_n\to\alpha$ であるから，ある自然数 $N_1$ が存在して，
	$n\geq N_1$ ならば
	\[
	|a_n-\alpha|<\frac{|\beta|\epsilon}{4}
	\]
	が成り立つ。
	
	また，$b_n\to\beta$ であるから，ある自然数 $N_2$ が存在して，
	$n\geq N_2$ ならば
	\[
	|b_n-\beta|
	<
	\frac{|\beta|^2\epsilon}{4(|\alpha|+1)}
	\]
	が成り立つ。
	
	ここで
	\[
	N=\max\{N_0,N_1,N_2\}
	\]
	とおく。
	
	このとき，$n\geq N$ ならば，
	\[
	n\geq N_0,\qquad n\geq N_1,\qquad n\geq N_2
	\]
	である。
	
	したがって，
	\[
	\frac{2}{|\beta|}|a_n-\alpha|
	<
	\frac{2}{|\beta|}\cdot \frac{|\beta|\epsilon}{4}
	=
	\frac{\epsilon}{2}
	\]
	である。
	
	また，
	\[
	\frac{2|\alpha|}{|\beta|^2}|b_n-\beta|
	<
	\frac{2|\alpha|}{|\beta|^2}
	\cdot
	\frac{|\beta|^2\epsilon}{4(|\alpha|+1)}
	\]
	であるから，
	\[
	\frac{2|\alpha|}{|\beta|^2}|b_n-\beta|
	<
	\frac{|\alpha|}{|\alpha|+1}\cdot\frac{\epsilon}{2}
	\leq
	\frac{\epsilon}{2}
	\]
	となる。
	
	よって，$n\geq N$ ならば
	\[
	\left|
	\frac{a_n}{b_n}-\frac{\alpha}{\beta}
	\right|
	\leq
	\frac{2}{|\beta|}|a_n-\alpha|
	+
	\frac{2|\alpha|}{|\beta|^2}|b_n-\beta|
	<
	\frac{\epsilon}{2}+\frac{\epsilon}{2}
	=
	\epsilon
	\]
	である。
	
	したがって，任意の $\epsilon>0$ に対して，ある自然数 $N$ が存在して，
	$n\geq N$ ならば
	\[
	\left|
	\frac{a_n}{b_n}-\frac{\alpha}{\beta}
	\right|<\epsilon
	\]
	が成り立つ。
	
	ゆえに
	\[
	\lim_{n\to\infty}\frac{a_n}{b_n}
	=
	\frac{\alpha}{\beta}
	\]
	である。
\end{proof}

\begin{proof}{B.2.1}
	まず，
	\[
	\lim_{n\to\infty}
	\sqrt{n}\left(\sqrt{n+1}-\sqrt{n}\right)
	=
	\frac{1}{2}
	\]
	を示す。
	
	左辺の数列を変形する。
	有理化すると，
	\[
	\sqrt{n}\left(\sqrt{n+1}-\sqrt{n}\right)
	=
	\sqrt{n}\cdot
	\frac{(n+1)-n}{\sqrt{n+1}+\sqrt{n}}
	\]
	である。したがって
	\[
	\sqrt{n}\left(\sqrt{n+1}-\sqrt{n}\right)
	=
	\frac{\sqrt{n}}{\sqrt{n+1}+\sqrt{n}}
	\]
	となる。
	
	分母分子を $\sqrt{n}$ で割ると，
	\[
	\frac{\sqrt{n}}{\sqrt{n+1}+\sqrt{n}}
	=
	\frac{1}{\sqrt{1+\frac{1}{n}}+1}
	\]
	である。
	
	よって，
	\[
	\sqrt{n}\left(\sqrt{n+1}-\sqrt{n}\right)
	=
	\frac{1}{\sqrt{1+\frac{1}{n}}+1}
	\]
	である。
	
	ここで，任意に $\epsilon>0$ を取る。
	示したいのは，十分大きな $n$ に対して
	\[
	\left|
	\frac{1}{\sqrt{1+\frac{1}{n}}+1}
	-
	\frac{1}{2}
	\right|
	<\epsilon
	\]
	が成り立つことである。
	
	実際，
	\[
	\left|
	\frac{1}{\sqrt{1+\frac{1}{n}}+1}
	-
	\frac{1}{2}
	\right|
	=
	\left|
	\frac{2-\left(\sqrt{1+\frac{1}{n}}+1\right)}
	{2\left(\sqrt{1+\frac{1}{n}}+1\right)}
	\right|
	\]
	であるから，
	\[
	=
	\frac{\sqrt{1+\frac{1}{n}}-1}
	{2\left(\sqrt{1+\frac{1}{n}}+1\right)}
	\]
	となる。
	
	さらに，
	\[
	\sqrt{1+\frac{1}{n}}-1
	=
	\frac{\frac{1}{n}}{\sqrt{1+\frac{1}{n}}+1}
	\]
	であるから，
	\[
	\left|
	\frac{1}{\sqrt{1+\frac{1}{n}}+1}
	-
	\frac{1}{2}
	\right|
	=
	\frac{1}
	{2n\left(\sqrt{1+\frac{1}{n}}+1\right)^2}
	\]
	を得る。
	
	ここで
	\[
	\sqrt{1+\frac{1}{n}}\geq 1
	\]
	であるから，
	\[
	\sqrt{1+\frac{1}{n}}+1\geq 2
	\]
	である。したがって
	\[
	\left(\sqrt{1+\frac{1}{n}}+1\right)^2\geq 4
	\]
	である。
	
	よって
	\[
	\left|
	\frac{1}{\sqrt{1+\frac{1}{n}}+1}
	-
	\frac{1}{2}
	\right|
	\leq
	\frac{1}{8n}
	\]
	が成り立つ。
	
	ここで，アルキメデスの性質より，
	\[
	N>\frac{1}{8\epsilon}
	\]
	を満たす自然数 $N$ を取ることができる。
	
	このとき，$n\geq N$ ならば
	\[
	\frac{1}{8n}
	\leq
	\frac{1}{8N}
	<
	\epsilon
	\]
	である。
	
	したがって，$n\geq N$ ならば
	\[
	\left|
	\sqrt{n}\left(\sqrt{n+1}-\sqrt{n}\right)
	-
	\frac{1}{2}
	\right|
	<
	\epsilon
	\]
	が成り立つ。
	
	ゆえに
	\[
	\lim_{n\to\infty}
	\sqrt{n}\left(\sqrt{n+1}-\sqrt{n}\right)
	=
	\frac{1}{2}
	\]
	である。
\end{proof}

\begin{proof}{B.2.2}
	次に，
	\[
	\lim_{n\to\infty}
	\frac{n^2+2n}{n+1}
	=
	\infty
	\]
	を示す。
	
	正の無限大への発散を示すには，任意の $M>0$ に対して，
	ある自然数 $N$ が存在して，$n\geq N$ ならば
	\[
	\frac{n^2+2n}{n+1}>M
	\]
	となることを示せばよい。
	
	まず式を変形する。
	\[
	\frac{n^2+2n}{n+1}
	=
	\frac{(n+1)^2-1}{n+1}
	\]
	であるから，
	\[
	\frac{n^2+2n}{n+1}
	=
	n+1-\frac{1}{n+1}
	\]
	となる。
	
	ここで，$n\geq 1$ のとき
	\[
	0<\frac{1}{n+1}<1
	\]
	である。したがって
	\[
	n+1-\frac{1}{n+1}>n
	\]
	が成り立つ。
	
	つまり，
	\[
	\frac{n^2+2n}{n+1}>n
	\]
	である。
	
	任意に $M>0$ を取る。
	アルキメデスの性質より，
	\[
	N>M
	\]
	を満たす自然数 $N$ を取ることができる。
	
	このとき，$n\geq N$ ならば
	\[
	n\geq N>M
	\]
	である。
	
	したがって
	\[
	\frac{n^2+2n}{n+1}>n\geq N>M
	\]
	となる。
	
	よって，任意の $M>0$ に対して，ある自然数 $N$ が存在して，
	$n\geq N$ ならば
	\[
	\frac{n^2+2n}{n+1}>M
	\]
	が成り立つ。
	
	したがって
	\[
	\lim_{n\to\infty}
	\frac{n^2+2n}{n+1}
	=
	\infty
	\]
	である。
\end{proof}